% TOP_PROBLEM: {"problemId": "problem.jones-unknot-conjecture", "canonicalTitle": "Jones Unknot Conjecture", "releaseRank": 339, "primaryDomain": "geometry_topology", "primaryDomainLabel": "Geometry & topology", "displayStatus": "open_with_solved_subcases", "releaseStatus": "open_with_solved_subcases"}
% !TeX program = lualatex
% ATTEMPT_STATUS: unresolved
\documentclass[11pt]{article}
\usepackage[margin=25mm]{geometry}
\usepackage{fontspec}
\setmainfont{Latin Modern Roman}
\usepackage{amsmath,amssymb,mathtools}
\usepackage{unicode-math}
\setmathfont{Latin Modern Math}
\usepackage{xurl,hyperref}
\hypersetup{colorlinks=true,urlcolor=blue}
\setcounter{secnumdepth}{0}
\setlength{\parindent}{0pt}
\setlength{\parskip}{5pt}
\setlength{\emergencystretch}{3em}
\begin{document}
\section*{\texorpdfstring{Jones Unknot Conjecture}{Jones Unknot Conjecture}}\label{339}
\subsection{Definitions and mathematical statement}
The Jones polynomial is the oriented-link invariant $V_L(t)$ normalized by $V_{\mathrm{unknot}}(t)=1$ and the skein relation
\[
 t^{-1}V_{L_+}(t)-tV_{L_-}(t)
 =(t^{1/2}-t^{-1/2})V_{L_0}(t),
\]
where the three diagrams differ only by a positive crossing, negative crossing, or oriented smoothing in a disk. For a knot $K\subset S^3$, the Jones unknot conjecture asserts
\[
 V_K(t)=1\quad\Longrightarrow\quad K\text{ is ambient isotopic to the unknot}.
\]
The equation is an identity of Laurent polynomials, not agreement at one numerical value of $t$. The converse follows from the normalization. The selected problem concerns knots, not a detection assertion for all links with several components.
\par\smallskip\textit{Formulation and concept references:} \href{https://encyclopediaofmath.org/wiki/Jones_unknotting_conjecture}{[S1]}.

\subsection{Short English statement}
Is the unknot the only knot whose entire Jones polynomial is the constant one?
\subsection{Research attempt: extreme states, adequate diagrams, and cancellation}
\textbf{Status and source scope.} The Jones unknot-detection assertion is not proved here. It concerns the entire normalized Laurent polynomial, rather than one value of it, and it concerns classical knots rather than arbitrary links. The primary computational paper of Tuzun and Sikora [E1] verifies the conjecture through twenty-four crossings. Its finite range is retained here; that computation was not rerun and is not a proof for unbounded crossing number. The source audit was performed on 27 September 2026.

The main partial result below is a direct state-sum proof of detection for diagrams that are both $A$- and $B$-adequate. This also identifies the exact cancellation obstacle in an attempted extension to arbitrary diagrams.

\textbf{Conventions for the state sum.} For a diagram $D$ with $c$ crossings, choose at each crossing one of its two Kauffman smoothings, called $A$ and $B$. A state $s$ has $a(s)$ choices of the first kind, $b(s)$ of the second, and $|s|$ resulting circles. Put
\[
 \delta=-A^2-A^{-2},\qquad
 \langle D\rangle=\sum_s A^{a(s)-b(s)}\delta^{|s|-1}.
 \tag{339.1}
\]
The empty crossing diagram consisting of one circle has bracket one. With the corresponding standard crossing convention, the normalized polynomial is
\[
 V_D(t)=(-A^3)^{-w(D)}\langle D\rangle,\qquad t=A^{-4},
 \tag{339.2}
\]
where $w(D)$ is the signed crossing number. The bracket relations give invariance under the second and third Reidemeister moves, and the factor in (339.2) compensates for the first move. They give the skein normalization in the statement. In particular multiplication by the writhe factor shifts all $A$-degrees equally and does not change their span.

Let $s_A$ and $s_B$ denote the numbers of circles in the all-$A$ and all-$B$ states. A diagram is $A$-adequate if, at every crossing, changing the all-$A$ smoothing there to $B$ merges two distinct state circles. Define $B$-adequacy by interchanging the roles. This is a condition on the diagram, not a condition imposed in the original question.

\textbf{Lemma 339.1 (no cancellation at adequate extremes).} For an $A$-adequate diagram the highest $A$-degree of its bracket is
\[
 c+2(s_A-1),
\]
and its coefficient is $(-1)^{s_A-1}$. For a $B$-adequate diagram the lowest degree is $-c-2(s_B-1)$, with coefficient $(-1)^{s_B-1}$.

\emph{Proof.} The all-$A$ contribution is $A^c\delta^{s_A-1}$ and has the asserted highest term. Consider another state with $r\geq1$ changes from the all-$A$ state. Make one of those changes first. Adequacy says it reduces the circle count by one. Every subsequent smoothing change alters the number of circles by either one or minus one. Thus the final state has at most $s_A-1+(r-1)=s_A+r-2$ circles. The highest degree in its contribution is at most
\[
 c-2r+2(s_A+r-3)=c+2s_A-6,
\]
which is four below the all-$A$ highest degree. No other state can cancel that term. The second assertion follows by starting with the all-$B$ state and reversing the degree inequalities. $\square$

\textbf{Corollary 339.2 (a genuine detection subcase).} If a knot has a diagram with $c>0$ which is both $A$- and $B$-adequate, its Jones polynomial is not one.

Indeed, the bracket span is
\[
 \operatorname{span}_A\langle D\rangle
   =2c+2(s_A+s_B-2)>0,
 \tag{339.3}
\]
since each state has at least one circle. The normalized polynomial has the same $A$-span and therefore cannot be constant. This is a complete argument within the stated diagram class.

For a connected reduced alternating diagram the checkerboard state circles give $s_A+s_B=c+2$, so (339.3) gives $\operatorname{span}_tV_D=c$. Here is the graph reason for the relation and the adequacy condition. Color the complementary regions of the connected projection black and white and form the plane checkerboard graph, with a vertex for each black region and an edge for each crossing. The alternating smoothing conventions identify the two state-circle collections with the vertex and face collections of this graph, in one order or the other. Euler's formula gives $v-c+f=2$. A loop in one state graph corresponds to a bridge in the dual graph, and in a reduced alternating diagram these would give a nugatory crossing. Their absence is exactly the required nonmerging failure excluded by adequacy. Thus the general adequate calculation recovers the alternating case with its sharper span formula.

\textbf{What goes wrong without adequacy.} If a crossing change of smoothing splits a single all-$A$ circle instead of merging two circles, a state with one $B$ choice can have $s_A+1$ circles. Its maximum possible degree is then
\[
 (c-2)+2s_A=c+2(s_A-1),
\]
the same as the all-$A$ contribution. The unique-extreme-term argument fails at that very step. Multiple states can now cancel at the apparent leading degree. Estimating only the degree of the all-$A$ and all-$B$ terms therefore cannot establish detection for an arbitrary diagram. This is a precise gap in the attempted generalization, rather than an appeal to computational difficulty.

\textbf{Connected sums provide another exact reduction.} State circles in a connected sum join one circle from each diagram, so $|s\#s'|=|s|+|s'|-1$. Equation (339.1) consequently factors, and writhe is additive. Therefore
\[
 V_{K\#L}(t)=V_K(t)V_L(t).
 \tag{339.4}
\]
If this product is one, each factor is a unit in the Laurent polynomial ring. We next check why a Jones polynomial which is a unit must itself be one, rather than an arbitrary monomial.

At $t=1$ the skein relation says that a crossing change does not change the value of the invariant. Crossing changes turn any knot into the unknot, so $V_K(1)=1$. For a two-component link, crossing changes turn it into the two-component unlink, giving $V_L(1)=-2$. The latter value follows either from the bracket or from the disjoint-circle factor $-(t^{1/2}+t^{-1/2})$.

Differentiate the skein relation for a crossing of a knot. Both $L_+$ and $L_-$ are knots, whereas the oriented smoothing $L_0$ has two components. At $t=1$ it gives
\[
 -1+V'_{L_+}(1)-1-V'_{L_-}(1)=V_{L_0}(1)=-2.
\]
Thus $V'_{L_+}(1)=V'_{L_-}(1)$. By crossing changes, $V'_K(1)=V'_U(1)=0$ for every knot. A unit of the Laurent polynomial ring is $\pm t^j$ (or $\pm t^{j/2}$ if the larger variable ring is used). Its value at one fixes the sign and its derivative there forces the exponent to be zero. Hence (339.4) yields
\[
 V_{K\#L}=1\quad\Longrightarrow\quad V_K=V_L=1.
 \tag{339.5}
\]
This rules out manufacturing a counterexample by multiplying two nonunit Jones polynomials and hoping for cancellation to the constant one. It does not prove that either knot factor with polynomial one is geometrically trivial.

\textbf{Exact computation and its complexity.} The state sum has $2^c$ states. Each can be evaluated with finite integer Laurent-polynomial arithmetic, and the number of resulting circles is at most $c+1$ for a connected knot projection. The bracket degrees lie between $-3c$ and $3c$, and the writhe normalization shifts them by at most $3c$. A coefficient is a sum of at most $2^c$ state contributions, each bounded in absolute value by $2^c$, so its bit length is $O(c)$. These bounds show that the final polynomial has a compact exact representation relative to the diagram size. They do not make its state-sum computation polynomial-time: the number of summands is still exponential.

In particular $V_K(1)=1$ cannot be used as a detector, because that equality holds for every knot. Exact interpolation from enough values would determine a degree-bounded Laurent polynomial, but computing those values at arbitrary points is another task. Replacing the full polynomial identity by a few experimentally chosen evaluations loses the asserted conclusion unless an adequate deterministic identity argument is supplied.

\textbf{Finite diagnostic.} The accompanying script evaluates the two-strand bracket algebra exactly. With $I$ the identity pairing and $e$ the cap-cup pairing, it uses $e^2=\delta e$ and a crossing $A I+A^{-1}e$. Closure sends $I$ to $\delta$ and $e$ to one. Odd powers give knot diagrams; the script checks the writhe normalization, the unknot case, and nonzero spans for the small nontrivial examples. This is a reproducible test of the normalization and algebra in (339.1)--(339.2), not a repetition of the twenty-four-crossing enumeration in [E1].

\textbf{Outcome and first gap.} Adequacy gives a complete extreme-term proof of detection, and connected sums reduce exactly as in (339.5). For arbitrary knots, the first missing ingredient is control of cancellation among states when the adequacy hypothesis fails. Neither the finite verification in [E1], the derivative identities at one, nor the polynomial-size output representation supplies that control. The selected conjecture remains unresolved by this attempt.

\subsection{Sources}
\begin{itemize}
\item[S1]Encyclopedia of Mathematics, "Jones unknotting conjecture"; conjecture arising from V.F.R. Jones, "Hecke algebra representations of braid groups and link polynomials," Ann. of Math. 126 (1987), 335-388.
\url{https://encyclopediaofmath.org/wiki/Jones_unknotting_conjecture}.
\textit{Formulation source.}
\item[E1]Robert E. Tuzun and Adam S. Sikora, \emph{Verification Of The Jones Unknot Conjecture Up To 24 Crossings}, arXiv:2003.06724v2 (2021).
\url{https://arxiv.org/abs/2003.06724v2}.
\textit{Primary computational source consulted 27 September 2026; its finite crossing bound is not extended here.}
\end{itemize}
\end{document}
